\documentclass[10pt,a4paper]{amsart}
\usepackage[margin=19mm]{geometry}
\usepackage{amsmath,amssymb,mathtools}
\usepackage[scr=rsfs,cal=euler]{mathalfa}
\usepackage{xfrac}
\usepackage[unicode,hidelinks,bookmarks=false]{hyperref}
\newcommand{\ie}{i.e.\ }
\newcommand{\cf}{cf.\ }
\newcommand{\resp}{respectively\ }
\newcommand{\Subsection}{\subsection}
\providecommand{\nobreakdash}{\nobreak\hbox{-}}
\setlength{\emergencystretch}{2em}
\multlinegap0pt
\usepackage{bm}
\usepackage{bbm}
\usepackage{dsfont}
\usepackage{xspace}
\usepackage{cancel}
\usepackage{mathtools}
\usepackage{amsthm}
\makeatletter
\@ifundefined{thm}{\newtheorem{thm}{Thm}}{}
\@ifundefined{lem}{\newtheorem{lem}[thm]{Lemma}}{}
\@ifundefined{thmA}{\newtheorem{thmA}{Theorem}}{}
\makeatother
\DeclareMathOperator{\Ball}{Ball}
\DeclareMathOperator{\dist}{d}
\DeclareMathOperator{\ran}{ran}
\DeclareMathOperator{\supp}{supp}
\newcommand{\cR}{\mathcal{R}}
\newcommand{\cU}{\mathcal{U}}
\newcommand{\ve}{\varepsilon}
\newcommand{\vs}{\varsigma}
\title[Contractibility of the automorphism group of a von Neumann a]{Contractibility of the automorphism group of a von Neumann algebra}
\author{Narutaka Ozawa}
\date{Source: arXiv:2412.17564v3 (CC BY 4.0)}
\begin{document}
\maketitle

\noindent\emph{Source and licence.} Contractibility of the automorphism group of a von Neumann algebra. Version arXiv:2412.17564v3 (CC BY 4.0), distributed under the Creative Commons Attribution 4.0 International licence, as recorded in the arXiv licence metadata for this exact version and shown on the abstract page https://arxiv.org/abs/2412.17564v3. Full rights notice: licenses/arxiv-2412.17564-CC-BY-4.0.txt.\par\smallskip

\noindent\textit{Editorial scope.} Counted unit: the Lem whose source label is \texttt{lem:pre} in the article (source0.tex lines 861--864, source label \texttt{lem:pre}), delivered together with the article's own complete original proof of it (source0.tex lines 865--964, one \texttt{proof} environment of 100 lines). No mathematics is written, repaired or re-derived: statement and proof are the article's own text line for line. Required context retained uncounted: thmB (lines 144--153); lem:selection3 (lines 813--827). Every definition, display and label of those lines is retained unchanged. Omitted from this package and not redistributed: 1--143, 154--812, 828--860, 965--1036. Nothing inside a retained excerpt is omitted or altered: every retained body is delivered exactly as the article's lines are written, so no omission receipt of any kind accompanies this package. Citations: the retained excerpts carry no citation command at all, so no bibliography entry of the article is transcribed and no reference is redistributed. Reference adaptation: none was needed and none was made; every reference inside the delivered unit resolves inside it, so no unresolved reference or citation is produced. Printed slips kept literal: no printed word, symbol, hypothesis or number of the counted statement or of its proof was altered. Layout and class: the article's own class is \texttt{amsart} with packages including mathtools; the wrapper is the lane's pinned \texttt{amsart} editorial wrapper at 10pt on A4, which restates the theorem-like environments the retained text needs and adds the article's own macro definitions that the retained text needs (\texttt{\textbackslash Ball}, \texttt{\textbackslash cR}, \texttt{\textbackslash cU}, \texttt{\textbackslash dist}, \texttt{\textbackslash ran}, \texttt{\textbackslash supp}, \texttt{\textbackslash ve}, \texttt{\textbackslash vs}), with extra wrapper packages (bm, bbm, dsfont, xspace, cancel, mathtools). The delivered numbering is the unit's own. Assets: no retained span contains a figure environment, a graphics-inclusion command, a PDF-page inclusion, a table or a TikZ picture; no image file is copied into this package, the package declares no asset dependency and no image file is redistributed with it. Licence: version arXiv:2412.17564v3 of this article is distributed under the Creative Commons Attribution 4.0 International licence, as recorded in the arXiv licence metadata for this exact version and shown on the abstract page https://arxiv.org/abs/2412.17564v3. A full rights notice is delivered at licenses/arxiv-2412.17564-CC-BY-4.0.txt. Characters: every retained excerpt is pure ASCII, so the delivered engine, XeLaTeX with its default Latin Modern text font, reports no missing character.
\begin{thmA}\label{thmB}
Let $N\subset M$ be an inclusion of $\sigma$-finite 
von Neumann algebras and assume 
that $N$ is strongly stable. 
Then the quotient map from $\mathcal{U}(M)$ 
onto $\mathcal{U}(M)/\mathcal{U}(N)$ 
admits a continuous cross section. 
Equivalently, there is a continuous equivariant retraction 
from $\cU(M)$ onto $\cU(N)$.
\end{thmA}

\begin{lem}\label{lem:selection3}
Let a continuous map $\sigma\colon X\to G$ 
and $\ve>0$ be such that 
$\dist(\sigma(x),Q^{-1}(x))<\ve$ for all $x$. 
Then for every continuous function $s\colon X\to(0,\infty)$ 
and $\delta>0$, 
there is a continuous function $t\colon X\to (0,\infty)$ 
that satisfies $0<t(x)\le s(x)$ and 
\[
Q^{-1}(x)\cap\Ball(\sigma(x),\ve)
 \cap \{ u\in G : \| u - \theta_{t(x)}(u) \|_2<\delta\}
 \neq \emptyset
\]
for every $x\in X$.
\end{lem}

\begin{lem}\label{lem:pre}
Theorem~\ref{thmB} holds true if the inclusion $N\subset M$ 
is strongly stable.
\end{lem}

\begin{proof} 
We stick to the notation and 
work with the $2$-norm 
corresponding to a faithful normal state 
$\omega=\omega_0\otimes\tau$ on $M=M_0\otimes\cR$ 
and the quotient map $Q\colon G\to X\coloneqq H\backslash G$.  
Put $\delta_n\coloneqq 4\cdot 2^{-n}$. 
We will construct continuous maps $\vs_n\colon X\to G$ 
and $s_n\colon X\to(0,\infty)$ such that 
\begin{itemize}
\item
$\vs_n(x)\in\ran\theta_{s_n(x)}$ for every $x\in X$;
\item
$\dist(\vs_n(x),Q^{-1}(x)) < \delta_n$ for every $x\in X$;
\item
$\dist(\vs_{n-1}(x),\vs_{n}(x)) < \delta_{n-1}+\delta_n$ for every $x\in X$.
\end{itemize}
Once this is done, $\vs(x) \coloneqq \lim_n \vs_n(x)$ defines 
a continuous cross section for $Q$.

Put $\vs_0(x)\coloneqq 1$ and $s_0(x)\coloneqq 1$ for all $x\in X$. 
Fix $n$ and suppose that $\vs_{n-1}$ and $s_{n-1}$ are already constructed. 
Put $\gamma\coloneqq 1/4$ and take $t\colon X\to (0,\infty)$ 
as in Lemma~\ref{lem:selection3} for 
$\sigma=\vs_{n-1}$, $\ve=\delta_{n-1}$, $s(x)=s_{n-1}(x)$, and $\delta=\gamma\delta_n$. 
For every $u\in G$, consider 
\begin{align*}
W_u\coloneqq\{ x \in X : &\ Q^{-1}(x) \cap \Ball(\vs_{n-1}(x), \delta_{n-1}) 
 \cap \Ball(u, \gamma\delta_n) \neq\emptyset\\
 &\mbox{ and } \| u - \theta_{t(x)}(u)\|_2 < \gamma\delta_n\}.
\end{align*}

We claim that $\{W_u\}_{u \in G}$ is an open cover for $X$. 
That it covers $X$ is straightforward. 
To prove $W_u$ is open, let $x_0\in W_u$ be given. 
Then there is $\kappa\in(0,\gamma\delta_n)$ such that 
\[
Q^{-1}(x_0) \cap \Ball(\vs_{n-1}(x_0), \delta_{n-1}-\kappa) 
 \cap \Ball(u, \gamma\delta_n) \neq \emptyset
\]
and $\| u - \theta_{t(x_0)}(u)\|_2 < \gamma\delta_n-\kappa$. 
Hence for 
\begin{align*}
&Q(\Ball(\vs_{n-1}(x_0), \delta_{n-1}-\kappa)
 \cap \Ball(u, \gamma\delta_n))\\
 &\quad\cap \{ x\in X : 
  \|\vs_{n-1}(x_0)-\vs_{n-1}(x)\|_2 + \| \theta_{t(x_0)}(u) - \theta_{t(x)}(u)\|_2<\kappa\}
\end{align*}
is an open neighborhood of $x_0$ that is contained in $W_u$. 

Thus, there is a partition of unity $\{\psi_i\}_{i\in I}$ for $X$ 
subordinated by $\{ W_u\}$. 
We may assume that $I$ is totally ordered. 
For each $i$, take $u_i$ such that $\supp\psi_i\subset W_{u_i}$. 
We define $p_i(x)$ to be the orthogonal projection in $L^\infty[0,1]$ that 
corresponds to the interval $[\sum_{j<i}\psi_j(x), \sum_{j\le i}\psi_j(x))$. 
Note that the defining sum is a locally finite sum and that 
the maps $x\mapsto p_i(x)$ are ultrastrongly continuous. 
We put $s_n(x)\coloneqq t(x)/2$ and define 
a continuous map $\vs_n\colon X\to G$ by 
\[
\vs_n(x) \coloneqq \sum_i \theta_{t(x)}(u_i)\phi_{t(x)}(p_i(x)) \in \ran\theta_{s_n(x)}.
\]
Recall that $\phi_t$ is a continuous family of 
trace-preserving embeddings of $L^\infty[0,1]$ 
into $\theta_t(\cR)'\cap\theta_{t/2}(\cR)$ 
(as embedded in $M$). 
To prove $\vs_n$ satisfies the displayed conditions, 
let $x\in X$ be given and write $I(x)\coloneqq\{ i : \psi_i(x)>0\}$. 
For each $i\in I(x)$, take 
\[
v_i \in Q^{-1}(x) \cap \Ball(\vs_{n-1}(x), \delta_{n-1}) \cap \Ball(u_i, \gamma\delta_n) 
\]
and put $v\coloneqq\sum_i \theta_{t(x)}(v_i)\phi_{t(x)}(p_i(x))$. 
Then, 
\[
\| v- \vs_n(x) \|_2\le \max\{\| v_i-u_i\|_2 : i\in I(x)\}< \gamma\delta_n.
\]
Moreover, for a fixed $i_0\in I(x)$, one has  
$v_iv_{i_0}^{-1} \in H$ for all $i$ and 
hence $v\theta_{t(x)}(v_{i_0})^{-1}\in H$, because 
$\theta_t$ and $\phi_t$ keep $N$ invariant. 
Recall that $\| u_{i_0} - \theta_{t(x)}(u_{i_0})\|_2 < \gamma\delta_n$ because 
$x\in W_{i_0}$.
Thus, 
\[
\dist( \vs_n(x), Q^{-1}(x) ) \le \| \vs_n(x) - v \|_2 + \| \theta_{t(x)}(v_{i_0})-v_{i_0}\|_2
\le 4\cdot \gamma\delta_n=\delta_n.
\]
Also, since 
$\|  \vs_{n-1}(x) - u_i\|_2 < \delta_{n-1}+\gamma\delta_n$ and 
$\vs_{n-1}(x) \in \ran\theta_{s_{n-1}(x)}\subset \ran\theta_{t(x)}$ 
commutes with $\phi_{t(x)}(p_i(x))$ for all $i\in I(x)$,
one has 
\[
\|\vs_{n-1}(x)-\vs_n(x)\|_2<\max\{ \| \vs_{n-1}(x) - \theta_{t(x)}(u_i) \|_2 : i\in I(x) \}
 < \delta_{n-1}+2\cdot\gamma\delta_n.
\]
These altogether verify the required properties for $\vs_n$ and $s_n$. 
\end{proof}

\end{document}
